\begin{figure}[t]
\centering\resizebox{.98\columnwidth}{!}{%
\begin{circuitikz}[american,x=1cm,y=1cm,font=\PLLfigurefont\small,
 c/.style={draw,minimum width=1.12cm,minimum height=.66cm,align=center},
 ar/.style={->,>=latex}]
 \node[c] (di) at (1.05,1.05) {TG}; \node[c] (ii) at (3.65,1.05) {INV};
 \node[c] (io) at (5.5,1.05) {INV}; \node[c] (fb) at (3.65,2.7) {TG};
 \draw (-.37,1.05) node[left]{$D$} -- (di.west);
 \draw (di.east) -- (ii.west); \draw (ii.east) -- (io.west);
 \draw (io.east) -- (6.67,1.05) node[right]{$Q$};
 \draw (6.32,1.05) node[circ]{} |- (fb.east); \draw (fb.west) -| (2.32,1.05) node[circ]{};
 \draw (di.north) -- ++(0,.28) node[above]{open};
 \draw (fb.north) -- ++(0,.16) node[above]{hold};
 \draw (4.58,1.05) node[circ]{} -- ++(0,.32) node[above]{$\overline Q$};
 \node[nmos] (mr) at (2.32,-.12) {};
 \draw (mr.D) -- (2.32,1.05); \draw (mr.G) -- ++(-.14,0) node[left]{RST};
 \draw (mr.S) -- ++(0,-.12) node[below]{VSS};
 \node[anchor=west,font=\PLLfigurefont\scriptsize] at (2.45,.6) {$X$};
 \node[anchor=west,font=\PLLfigurefont\scriptsize] at (2.65,-.1) {$W_n=2~\mu$m};
 \begin{scope}[shift={(0,-2.55)}]
 \node[c] (ml) at (1.5,0) {Master latch};
 \node[c] (sl) at (4.8,0) {Slave latch};
 \draw[ar] (-.35,0) node[left]{$D$} -- (ml.west);
 \draw[ar] (ml.east) -- node[above]{$Q_M$} (sl.west);
 \draw[ar] (sl.east) -- (6.65,0) node[right]{$Q$};
 \node[below,align=center] at (ml.south) {open $L$\\hold $H$};
 \node[below,align=center] at (sl.south) {open $H$\\hold $L$};
 \end{scope}
\end{circuitikz}}
\caption{R1 resettable static latch (top) and ordinary positive-edge DFF (bottom). Each TG receives an enable and its complement. The feedback TG returns $Q$ to storage node $X$ during hold. Active reset pulls $X$ to VSS in each latch.}
\label{fig:r1afc-latch}
\end{figure}

Two latches form one DFF. The preset variant inverts $D$ and exposes
the slave's complementary output. Each shared clock group forms $K_d$
with six cascaded inverters. With active-high reset $R$, the interlock
computes
\begin{align}
 L &= \overline K\,\overline{K_d}\,\overline H\,\overline R, &
 H &= K\,K_d\,\overline L\,\overline R .
\end{align}
A NAND4 and two large inverters drive each phase and its complement.

R1 contains 146 DFFs and one additional latch on the feedback clock
gate. The 293 latches use 586 TGs, 586 latch inverters and 293 reset
NMOS devices. Outside these latches, the generator emits 531 inverters
(including clock drivers), 740 NAND2s, 79 NOR2s and 20 NAND4s. The
total is $293(9)+531(2)+740(4)+79(4)+20(8)=7{,}135$ MOS instances
in the AFC core. These counts can be checked against the selected R1
gate generator, its port manifest and the generated transistor netlist.
